\chapter{绪论}
本章简要介绍本项目的研究背景、目标以及主要内容。项目旨在开发一个高效的 Search Agent 系统，该系统通过连接 MCP server 和后端搜索引擎，实现对用户查询的智能处理与精准搜索。系统包括两大核心功能模块：实时在线搜索和离线搜索。在线搜索模块支持基于实时数据的动态爬取与内容展示，离线搜索则专注于垂直领域的数据处理与检索。通过关键词提炼、文本检索与向量检索相结合，系统能够有效提升搜索精度与响应速度。
\section{研究背景与意义}
 
随着互联网内容规模和多样性的快速扩展，用户对信息获取的即时性、准确性和场景相关性的要求不断提高。在这个大数据时代，信息的生成速度与更新频率已经远超传统检索系统的处理能力。特别是用户从搜索引擎、社交媒体、新闻网站等平台获取信息的需求日益增长，传统基于关键字的检索方法在面对高维语义查询、非结构化内容和跨媒体信息时逐渐显露出局限性\cite{mao2025comparativestudyretrievalmethods}。这些方法主要依赖词汇匹配或字符串匹配技术，尽管在某些场景下能有效检索到符合条件的结果，但它们对语义理解的能力较弱。尤其在需要捕捉用户深层意图和跨领域知识的情境中，基于关键字的检索方法往往无法提供满意的检索效果。此外，这些方法对于处理复杂查询、长尾问题以及多模态数据时显得力不从心，导致信息覆盖面不足和用户体验不佳。

在此背景下，语义向量检索、预训练大模型以及混合检索策略应运而生，并为提升检索覆盖与相关性提供了新的可能。语义向量检索通过将文本、图像等数据转化为向量空间中的点，能够更好地捕捉到信息之间的语义关联，突破了传统基于关键字的匹配限制。通过使用深度学习技术训练的大型预训练模型，系统可以实现更为精确的语义理解，从而大幅提高信息检索的质量与效率\cite{wang2025balancingblendexperimentalanalysis}。混合检索策略则结合了传统检索与基于语义向量的检索方法，利用两者的优势，进一步增强了检索结果的精准度与覆盖面，尤其在大规模数据环境中，能够显著优化用户的检索体验\cite{yang2025arcaderealtimedatahybrid}。

与此同时，商业与工程实践中存在两类互补的数据来源，这对现代检索系统的架构设计提出了更高的要求。第一类是需要保证时效性和新鲜度的实时在线内容，这类内容的更新频繁且与时事新闻、社交动态等紧密相关，因此其检索要求极高的响应速度与及时更新能力\cite{ALAKASHI2022}。第二类是具有专业垂域价值的离线数据集合，这类数据一般来源于学术论文、行业报告、专业文献等，包含了大量的结构化或半结构化的信息。虽然这些数据的时效性较差，但其内容的深度和专业性使得它们在某些特定场景下具有不可替代的价值。

然而，实时在线数据和离线垂域数据在检索需求、索引策略与召回机制上存在显著差异。实时数据要求系统具有高效的抓取与更新能力，以应对内容的快速变动，并且在时间限制下提供准确的结果；而离线数据则更侧重于高质量的索引与深度检索，系统需要在保证召回覆盖的同时，尽可能提高检索效率。这就要求检索系统能够同时应对两类数据的特点，综合考虑时效性与专业性，提出合适的检索策略。


本研究提出并实现以 Search Agent 为核心的工程化方案，旨在通过模型驱动的 prompt 解析将自然语言提示稳定、准确地转化为检索指令\cite{10.1145/3748304}，进而在 MCP 协调下高效触发实时在线检索或离线检索通道\cite{luo2025evaluationreportmcpservers}，并在必要时发起实时爬取以增强结果的新鲜度与信息深度。该工作具有重要工程与学术意义：一方面，从工程视角，它覆盖了从提示理解、检索调度、实时爬取、实时解析、到数据清洗与入库的端到端流程，解决了多通道协同、时延控制与容错调度等实用问题，为生产环境下的检索平台提供可落地的实现路径；另一方面，从方法论角度，结合模型驱动的语义解析与混合检索策略\cite{kalra2025morbetterhandlingdiverse}，可以显著提升召回覆盖与排序质量，探索在低延迟约束下的召回融合与分层重排策略，对后续检索系统的研究与优化具有借鉴意义。因此，构建一个工程化且可部署的模型驱动实时—离线混合检索 Search Agent，不仅能满足实际业务需求，也具有推动检索系统向更语义化、更实时化方向发展的学术价值。

\section{研究发展现状} \subsection{信息检索与混合检索的发展} 信息检索领域自经典倒排索引与布尔检索起步，经历了概率检索模型、基于语言模型的检索方法，到近年来的基于深度学习与表示学习的语义检索体系演进\cite{raj2025exploringnewapproachesinformation}。倒排索引在处理大规模文本检索时以其高效的召回与成熟的工程实现成为主力，而向量检索（尤其是基于深度语义表示的近似最近邻搜索）在语义匹配与模糊查询场景下展现出显著优势\cite{xu2026surveymodelarchitecturesinformation}。由此催生的混合检索架构，通常通过并行或分层方式同时利用倒排索引与向量索引来获得既有精确匹配又有语义覆盖的候选集。近年的研究与工程实践还重点关注召回融合策略，如分数级融合、RRF等算法已被广泛应用，而基于原始分数的动态加权融合策略可进一步提升；索引更新与增量构建策略以降低索引维护成本并保证数据新鲜度，其中流式向量量化与双索引架构通过分离操作更新与检索查询来实现近实时的索引变更可见性\cite{prajapati2025livevectorlakerealtimeversionedknowledge}；以及在有限资源下的检索效率优化，例如基于乘积量化的近似最近邻搜索压缩技术，通过聚类划分空间与量化大幅降低存储开销，以及通过即时重算嵌入而非存储的存储高效型索引方案\cite{wang2025leannlowstoragevectorindex}。

在工程化方向，检索系统趋向模块化与服务化，出现了分离的索引服务、检索召回服务、重排服务与上下游数据流水线\cite{10.1145/3726302.3730351}。实时检索场景推动了流式索引、近实时更新与低延迟 API 的发展，而离线垂域数据的采用则强调批处理编目、质量筛选与语义向量化的一致性。总体来看，学术界在召回模型、向量表示与重排模型上持续推进，工业界则在系统可用性、可扩展性与运维方面积累了大量经验，混合检索正逐步由实验室成果向生产系统演进。

\subsection{大模型在提示解析与检索驱动中的作用} 
预训练大模型在自然语言理解、生成及抽取任务上展示出强大的通用能力，使其在将用户自然语言提示转化为检索动作中具备天然优势。通过对 prompt 进行语义分词、意图识别、实体抽取与关键词提炼，模型能够为检索器提供更符合用户语义需求的检索词或扩展查询，从而改进召回的相关性。在工程实践中，模型驱动的 prompt-to-query 模块\cite{li2023agent4rankingsemanticrobustranking}通常需要具备稳定性、可解释性与可控性，这意味着在使用大模型时应结合规则化后处理、置信度评估与模板约束以降低噪声引入风险\cite{huang2026qponemodelunifiedgenerativellm}。

此外，大模型可用于对检索结果的初步筛选或摘要生成，为是否触发下一步实时爬取或展示决策提供支持。模型还能够在质量评估阶段对爬取内容进行语义一致性检查与关键信息抽取\cite{zhang-etal-2025-structured}，从而提升二次清洗的自动化水平。然而，将大模型引入实时系统也带来了挑战，包括推理延迟、资源消耗与模型输出的可靠性问题\cite{shachar2025nerretrieverzeroshotnamed}。为此，实际工程常采用混合策略：在延迟敏感路径上使用轻量化或蒸馏模型，而在非实时或离线环节使用更强的模型以保证质量；同时引入缓存、异步化与并行化机制以缓解延迟瓶颈。

\subsection{实时检索与离线垂域检索的工程需求与应用场景} 
实时检索强调低延迟与高新鲜度\cite{li2023realtime}，应用于热点事件监测、即时问答与短时话题聚合等场景。工程上需要解决的问题包括实时索引或近实时索引更新\cite{chen2024near}、并发与超时控制\cite{wang2023concurrency}、快速判定是否需要发起深度爬取\cite{zhang2024deep}以及保证爬取与解析在可接受延迟内完成\cite{zhao2023lowlatency}。实时爬取除了常规的反爬、限速与容错外，更要求解析模块能处理流式数据并将结构化信息及时下发至后续清洗与入库流程，从而缩短从采集到可检索之间的时延。此外，实时检索系统还需处理并发与超时控制、快速判定是否需要发起深度爬取以及保证爬取与解析在可接受延迟内完成。实时爬取除了常规的反爬、限速与容错外，更要求解析模块能处理流式数据并将结构化信息及时下发至后续清洗与入库流程，从而缩短从采集到可检索之间的时延。在爬取侧，基于流式处理框架构建的低延迟可扩展爬虫可将 URL 抓取与解析作为持续流处理\cite{stormcrawler}；在实时数据摄入侧，基于 CDC 的变更数据捕获管道可将源数据库的增、删、改操作实时流式同步至向量数据库，避免批量更新的滞后问题。多模态实时数据系统也展示了跨文本、向量、空间等不同数据模态的连续查询与高吞吐摄入能力。

离线垂域检索则更多面向内容丰富、具有较高语义价值或社区特征的数据源，如小红书、知乎等。离线通道的关键在于高质量索引构建、向量化一致性、批量数据清洗与标签化，并保证在检索时能够提供领域内高度相关的候选项\cite{sun2023offline}。离线数据通常需要更复杂的质量评估机制、去噪与样式适配，以及周期性更新策略以兼顾新旧数据的平衡\cite{wu2023periodic}。现实应用中，实时与离线通道往往需要协同工作：实时通道补充新鲜度，离线通道提供深度与专业性。因此工程上需设计统一的调度与融合机制，使两类通道在延迟预算内实现最优召回与排序效果\cite{yang2026routirfastservingretrieval}，同时保证系统的可运维性与扩展性。
\subsection{查询算法的分类和发展现状}

查询算法作为信息检索领域的核心技术，旨在通过高效的搜索策略从海量数据中获取与用户查询最相关的结果。随着大数据和人工智能技术的飞速发展，查询算法的研究与应用也经历了显著的变革。尤其是在倒排、向量和重排等技术的结合与发展下，查询算法的性能得到了显著提升。本文将重点回顾倒排、向量和重排技术的研究进展，并分析它们在查询算法中的应用现状与未来趋势。

倒排技术是传统信息检索中最为基础且关键的组成部分之一。倒排索引通过建立词项到文档的映射关系，大大提高了查询过程中的效率。在倒排索引的构建中，首先会对所有文档进行分词、去除停用词并记录词频，从而生成一个倒排表格，方便快速查询。该技术的优势在于其结构的紧凑性和查询效率，特别是在处理大规模数据集时，倒排索引能够显著减少搜索时间。虽然倒排技术在早期已取得较大成功，但随着数据规模的不断扩大，传统的倒排索引方法在处理复杂查询、多语种信息和多模态数据时逐渐显现出局限性。因此，近年来，研究者开始探索更为高效的倒排结构，例如基于压缩技术的倒排索引，以及多层次、分布式的倒排索引方法，旨在进一步提高查询效率和存储优化。

与倒排索引相辅相成的是向量空间模型，该模型通过将文本表示为向量形式，进而计算文档与查询之间的相似度。向量空间模型在信息检索中扮演着重要角色，尤其是在查询结果排序和相关性度量方面。在该模型中，每个文档被表示为一个高维向量，向量的维度对应于文档中的词汇。通过计算查询向量与文档向量之间的相似度，查询算法能够筛选出与用户查询最为相关的文档。然而，传统的向量空间模型也面临着数据稀疏性、向量维度过高等问题，导致查询效率和效果的下降。为了解决这些问题，近年来，研究者提出了基于降维、主题模型以及词嵌入技术的向量表示方法，这些方法能够在一定程度上缓解数据稀疏性问题并提升查询精度。

在倒排索引和向量模型的基础上，倒排向量重排技术逐渐成为提升查询精度的重要研究方向。倒排向量重排旨在通过对初步查询结果进行重新排序，进一步提高检索结果的相关性。在传统的信息检索中，倒排索引能够快速检索出相关文档，但往往无法准确区分文档之间的细微差异，导致检索结果的准确性不足。为此，研究者引入了重排机制，结合机器学习、深度学习等方法，通过学习文档和查询之间的复杂关系，调整倒排结果的排序。最初的倒排重排方法多依赖于传统的排序算法，如支持向量机、决策树等，这些方法通过手工设计特征进行训练，尽管在某些场景下取得了较好效果，但其对特征工程的依赖较大，且对复杂查询的处理能力有限。

近年来，深度学习技术的引入为倒排向量重排带来了新的突破。基于深度神经网络的倒排重排方法能够自动学习更加复杂的特征，避免了传统方法中繁琐的特征工程。例如，基于卷积神经网络和循环神经网络的模型，能够捕捉到文档与查询之间的语义关系，提高重排的准确性。此外，强化学习的应用也为倒排重排提供了新的思路，模型通过与环境的交互不断调整重排策略，优化查询结果的质量。基于深度学习和强化学习的倒排向量重排技术，在多个实际应用中，如搜索引擎、推荐系统和问答系统等，展现出了优越的性能。
\section{研究目标与任务} 
本研究的总体目标是设计并实现一个以模型驱动的 Search Agent 为核心、能够在 MCP 协调下对接实时在线与离线检索后端的工程化检索流水线，满足低延迟响应、高检索质量与可部署运维的要求。MCP 通过统一的客户端—服务器协议实现了 AI 代理与多种数据源之间的通信标准化，为检索系统提供了可扩展的架构基础\cite{sun2025symbioticragenhancingdocumentintelligence}。在 MCP 生态中，ETOM 评估框架进一步对多跳端到端工具协调能力进行系统诊断，发现检索错误占所有失败原因的近一半，凸显了检索环节在 MCP 协调架构中的瓶颈地位\cite{Dong2025ETOMAF}。为实现该目标，论文将完成以下具体任务：

首先，构建 prompt-to-query 模块，采用预训练模型结合规则化后处理方法实现对用户自然语言提示的语义分词、关键词提炼与检索意图判别，并给出置信度评估与可控的输出格式以便于下游调用。在意图检测方面，SAID 框架通过同时建模文本信息与查询—答案关系结构，在小样本场景下显著提升了查询意图检测的准确性。在查询改写方面，AdaRewriter 采用测试时自适应策略，利用轻量级奖励模型在推理阶段从候选改写中筛选最优查询，在多个对话式检索数据集上取得了显著优于现有方法的性能\cite{lai2025adarewriterunleashingpowerpromptingbased}。

其次，设计实时在线检索通道的时序与控制策略，包含热点定时检索、超时与并发控制、检索结果的实时下发机制，以及在关键词指示需要详细展示时触发的 Python 实时爬取模块，爬取模块需包含动态限速、超时处理、实时解析与断点容错能力。在实时数据采集架构方面，FediLive 框架面向去中心化社交网络构建了高并发、高容错、易扩展的跨实例爬取与数据预处理系统，通过动态限速、轮询去重、ID 冲突解决及多线程调度等技术，在有限计算资源下实现了大规模实例的全平台采集，为实时爬取模块的设计提供了可复用的工程方案。

第三，构建离线检索通道的索引与向量化管线，完成离线数据的采集、批量清洗、向量化建库与倒排索引的增量更新策略。在数据清洗环节，LOCA 框架通过“增强—审阅”循环自动补全缺失的逻辑步骤并区分科学原理与推导过程，能够将科学语料中的错误率从 20\% 降低至 2\% 以下，展示了大语言模型在自动化数据清洗与质量评估中的潜力\cite{fang2025localogicalchainaugmentation}。在向量索引的增量更新方面，LiveVectorLake 提出了双层时间知识库架构，将热层向量索引与冷层列式版本存储分离，实现了当前知识的毫秒级检索与历史版本的点时间追溯。

第四，搭建端到端的数据处理流水线，包括实时解析、二次清洗、质量评估、ETL 与入库，并实现基础监控、日志与故障恢复机制以保证系统的工程可用性。

通过上述任务的实施，最终在吞吐与延迟、检索与排序质量、新鲜度以及工程指标上验证系统的实用性与可部署性。

\section{论文结构安排} 
本文围绕Search Agent系统的设计与实现展开，整体结构共分为六章，逻辑层层递进，由理论综述逐步过渡至系统设计、实现与验证。

第一章为绪论，主要介绍研究背景与意义，结合当前信息检索技术的发展趋势，分析混合检索与大语言模型结合的研究价值。同时，对国内外研究现状进行梳理，重点涵盖信息检索技术演进、大模型在提示解析与检索驱动中的作用，以及实时检索与离线垂域检索的工程需求。在此基础上明确论文的研究目标与主要任务，并对全文结构进行总体说明。

第二章为技术综述，系统梳理本文所涉及的关键技术与工具，包括开发与性能工具、数据监控体系以及混合检索中的核心算法，如Lucene倒排索引、Faiss向量检索及学习排序模型等，为后续系统设计提供技术基础。

第三章为系统需求分析与总体设计，是全文核心部分之一。首先从系统功能与业务流程出发，对Prompt解析、实时检索、离线检索、索引构建等模块进行需求分析；随后从逻辑视图、过程视图、开发视图与物理视图四个层面进行系统架构设计，并详细说明各模块划分及其交互关系，同时完成数据库与索引存储方案设计。

第四章为系统实现，具体阐述各核心模块的实现方法，包括Prompt解析、实时检索与爬取、离线检索与索引构建以及召回融合与重排机制，重点体现系统工程落地过程与关键技术细节。

第五章为系统测试，从单元测试、功能测试及非功能测试三个层面，对系统的正确性、稳定性与性能进行验证，确保系统满足设计预期。

第六章为总结与展望，对全文工作进行归纳，并对未来可能的优化方向与研究拓展进行讨论。